\section{Confronto finale tra FairMind e LLM: tempo, token e accuratezza}
\label{sec:confronto-finale}

L'esecuzione automatica del notebook \texttt{2\_3\_benchmark\_thor.ipynb} tramite \texttt{sbatch run\_benchmark.sbatch} (job Slurm 45149, 20 luglio 2026, dopo la correzione descritta in \S\ref{sec:binning} e la risoluzione dei problemi di infrastruttura descritti in \S\ref{sec:thor}) costituisce il primo confronto end-to-end completamente riproducibile su Thor, e produce i seguenti risultati:

\begin{center}
\begin{tabular}{lrrr}
\toprule
Effetto & FairMind & LLM & Errore relativo \\
\midrule
TV & 0.194470 & 0.194500 & 0.02\% \\
TE & 0.183161 & 0.231167 & 26.21\% \\
SE & -0.007296 & -0.036667 & 402.56\% \\
DE & 0.137359 & 0.083167 & 39.45\% \\
IE & -0.045803 & -0.047000 & 2.61\% \\
\bottomrule
\end{tabular}
\end{center}

FairMind completa il calcolo in circa 8 millisecondi; il LLM impiega 27.42 secondi (44 secondi il job Slurm nel suo complesso, includendo l'avvio del container), consumando 6870 token in ingresso e 2037 in uscita (8907 totali) --- circa 3500 volte più lento di FairMind, a fronte di un prompt di sole cinque tabelle pre-aggregate anziché dati grezzi.

Confrontando questo risultato con l'intera evoluzione documentata nelle sezioni precedenti emerge una traiettoria coerente: dal primo esperimento con o3-mini su 300 righe (265 secondi, oltre 102\,000 token totali, \S\ref{sec:analisi-json}), agli esperimenti sui livelli di \emph{reasoning effort} (192--203 secondi, fino a 85\,021 token, senza guadagni di accuratezza consistenti, \S\ref{sec:effort}), al tentativo con 2000 righe grezze su Thor (fino a circa 32\,000 token e risultati instabili, \S\ref{sec:limiti}), fino all'approccio a tabelle pre-aggregate qui utilizzato (8907 token, risultato stabile e riproducibile). Il costo per interrogazione si è ridotto di oltre un ordine di grandezza lungo il percorso, mentre l'accuratezza sugli effetti più semplici (TV, IE) è rimasta buona fin dai primi esperimenti ed è oggi la più solida del confronto finale.

Va segnalato un punto che il repository non permette di chiarire con certezza: il codice (\texttt{src/llm.py}, campo \texttt{model="qwen2.5-7b-instruct"}) dichiara l'uso di Qwen2.5-7B, ma lo script Slurm che avvia il server persistente effettivamente interrogato in questo esperimento (\texttt{slurm/llama\_server\_gpu.sbatch}, verificato il 20 luglio direttamente sul cluster) carica il file \texttt{Qwen2.5-14B-Instruct-Q4\_K\_M.gguf}. Poiché il server \texttt{llama.cpp} ne serve uno solo indipendentemente dal nome richiesto nella chiamata API, è plausibile che i risultati qui riportati --- e verosimilmente anche quelli storici raccolti nelle sessioni precedenti sullo stesso server --- provengano dal modello da 14 miliardi di parametri e non dal 7B sempre dichiarato in codice e documentazione. Si tratta di un'ipotesi supportata da evidenza indiretta, non di un fatto verificato con certezza assoluta, e andrebbe confermata prima di riportare il numero di parametri del modello nella versione finale della tesi.
